\documentclass[11pt]{article}
\usepackage[a4paper,margin=1in]{geometry}
\usepackage{amsmath,amssymb}
\usepackage{hyperref}
\title{Orientation Space, Symmetry Reduction, and Disorientation in PyTex}
\author{PyTex Contributors}
\date{\today}

\begin{document}
\maketitle

\section{Purpose}

PyTex represents orientations by unit quaternions and rotation matrices while keeping crystal and specimen reference frames explicit. The central Phase 2 question is how these rotations are reduced under symmetry.

\section{Orientation Equivalence}

If $g$ is an orientation mapping crystal coordinates into specimen coordinates, then symmetry-equivalent orientations take the form
\[
  g' = s g c,
\]
where $s$ belongs to specimen symmetry and $c$ belongs to crystal symmetry.

\section{Disorientation}

Given two orientations, PyTex computes a base misorientation and then searches the symmetry orbit for the minimum-angle representative. This minimum-angle representative is treated as the current disorientation used in kernel evaluation and neighborhood calculations.

\section{Exact Orientation Representative}

For stable orientation-space comparison, PyTex now exposes an exact orbit-reduction rule in the quaternion hemisphere. For each symmetry-equivalent orientation candidate, it computes a canonical quaternion with non-negative scalar part and then selects the representative with maximal scalar part, equivalently minimal rotation angle to the identity.

This gives an exact minimum-angle representative even when no workflow-specific reference orientation is provided, with a deterministic tie-break for symmetry-boundary cases.

\section{Current Implementation Boundary}

PyTex presently implements exact orbit reduction for the supported proper point groups already present in the codebase and minimum-angle disorientation. What remains ahead is a proof-oriented class-by-class boundary catalog and broader external parity on those exact boundaries.

\section{Normative References}

\begin{thebibliography}{9}
\bibitem{bunge}
H.-J. Bunge,
\textit{Texture Analysis in Materials Science: Mathematical Methods},
Butterworths, 1969.
DOI: \url{https://doi.org/10.1016/C2013-0-11769-2}.

\bibitem{degraef}
M. De Graef,
\textit{Introduction to Conventional Transmission Electron Microscopy},
Cambridge University Press, 2003.
DOI: \url{https://doi.org/10.1017/CBO9780511615092}.
\end{thebibliography}

\end{document}
